\documentclass{article}
\usepackage{graphicx}
\usepackage{amsmath}
\usepackage{amssymb}
\usepackage{amsfonts}
\usepackage{hyperref}
\usepackage{exercises}
\usepackage{indentfirst}
\usepackage{circledsteps}
\usepackage{amsthm}
\usepackage{csquotes}
\usepackage{yhmath}
\usepackage{comment}
\usepackage{placeins}
\usepackage{geometry}
\usepackage{hyperref}
\usepackage{tikz}
\newcommand*{\oldneg}{\mathord{\sim}}
\newcommand{\Z}{\mathbb{Z}}
\newcommand{\N}{\mathbb{N}}
\let\oldemptyset\emptyset

\theoremstyle{definition}
\newtheorem{definition}{Definition}
\newtheorem{theorem}{Theorem}[section]
\newtheorem{corollary}{Corollary}[theorem]
\newtheorem{lemma}[theorem]{Lemma}

\setlength{\parindent}{.5in}
\geometry{margin=0.5in}
\setlength{\parskip}{12pt}
\title{Chapter 10 Notes}
\begin{document}
\date{}
\maketitle

\section*{Section 1}

\subsection*{Exercise 1.1}
Suppose that $(A,R)$ is isomoprhic to $(A',R')$ under $\phi$.
In fact, the biconditional does not need to be proven for any of these - proving the
implication also proves the converse.

(a) Suppose $(A,R)$ is proper. Then, for every $x\in A$, $x^*$ is a set. Now consider any 
$x\in A'$. If $y\in x^*$, then $\phi^{-1}(y)\in \phi^{-1}(x)^*$. And, if 
$y\in\phi^{-1}(x)^*$, then $\phi(y)\in x^*$. Thus there is a 1-1 correspondence between 
$x^*$ and $\phi^{-1}(x)^*$, and because $\phi^{-1}(x)^*$ is a set, so is 
$x^*$ (by the axiom of substitution).

(b) Suppose $(A,R)$ is extensional. Then, for all elements $x,y$ in $A$, if 
$x\ne y$, then there is at least one element $z$ in $A$ such that $z$ is a component of 
one of $x$ or $y$, but not of both. Consider any elements $x,y$ in $A'$. If $x\ne y$, 
then $\phi^{-1}(x)\ne\phi^{-1}(y)$, meaning, because $(A,R)$ is extensional, there is some 
$z\in A$ such that $z$ is a component of one of $\phi^{-1}(x)$ or $\phi^{-1}(y)$, but not of 
both. But, because $(A,R)$ and $(A',R')$ are isomorphic, it follows that $\phi(z)$ is a component of 
$x$ or $y$, but not both, meaning $(A',R')$ is extensional.

(c) Suppose $(A,R)$ is well founded. Thus every subclass of $A$ has an initial element. 
Consider any subclass $B'$ of $A'$. Then the class $B$ of all 
$x\in A$ such that $\phi(x)\in B'$ is in a 1-1 correspondence to $A'$. Then 
$B$ has some initial element $x$. Then $\phi(x)\in B'$. If, for any 
$y\in\phi(x)^*$, $y\in B'$, we would have $\phi^{-1}(y)\in B$. But 
$\phi^{-1}(y)\in x^*$, contradicting the fact that $x$ is an intial element of 
$B$. Hence $\phi(x)^*\cap B=\emptyset$, meaning $(A',R')$ is extensional.

\end{document}
